\begin{claim}{B.3}{Proved}{Relative Theorem A for an arbitrary derivator}
Let $\D$ satisfy Der~1--4 in the presheaf convention of A.1, on all
$\U$-small categories. Let $A\xrightarrow{u}B\xrightarrow{q}S$ be
functors of $\U$-small categories and put $p=qu$. If, for every
$s\in S$, the specified restriction
\[
 (u/s)^*:\LC_\D(B/s)\longrightarrow\LC_\D(A/s)
 \tag{B.3.1}
\]
is an equivalence, then $u^*:\LC_\D(B)\to\LC_\D(A)$ is an
equivalence. If every restriction in (B.3.1) is fully faithful,
then $u^*$ is fully faithful. Both conclusions also hold for a map
of Grothendieck opfibrations preserving cocartesian arrows when the
corresponding restriction on every strict fiber is, respectively,
an equivalence or fully faithful. Consequently $\Wlc$ is a
fundamental localizer for every such derivator, while $\Wone$
satisfies relative Theorem A despite its failure of
two-out-of-three in A.5. The statement concerns the ordinary categories
$\LC_\D(-)$ themselves; neither an infinity-categorical presentation
nor strongness of $\D$ is required.
\end{claim}
\begin{proof}
We first establish the partial localizations needed to make descent
coherent. If $T$ is small and $R\subseteq T$ is a subcategory, let
$\D(T;R)$ be the full subcategory of $\D(T)$ on objects whose
underlying presheaf inverts every arrow of $R$. Both adjoints to
\[
 j_{T,R}:\D(T;R)\hookrightarrow\D(T)
 \tag{B.3.2}
\]
exist, and they are the components of adjunctions of shifted
derivators. Here and below ``shifted'' means that, at an auxiliary
shape $J$, the same condition is imposed on the $R$ arrows in
$T\times J$.

For completeness, this is the arbitrary-subcategory form of the
argument in \citet[Theorem~3.6, Lemma~3.9, Proposition~3.10, and
Theorem~3.13]{Cisinski2009LocallyConstant}. Write $C=T^{\op}$ and
$R'=R^{\op}$. The universal derivator of simplicial presheaves on
$C$ has a localization
\[
 \gamma:\Ho([C^{\op},\mathrm{sSet}])
 \rightleftarrows
 \Ho(L_{R'}[C^{\op},\mathrm{sSet}]):i,
 \qquad\gamma\dashv i.
 \tag{B.3.3}
\]
The right adjoint $i$ is fully faithful. It preserves homotopy
colimits: a pointwise homotopy colimit of diagrams which invert
$R'$ again inverts $R'$, since a homotopy colimit sends an
objectwise weak equivalence of diagrams of spaces to a weak
equivalence. This argument works at every auxiliary shape. Thus both
$\gamma$ and $i$ are cocontinuous morphisms of derivators.
Cisinski's Theorem~3.6, which is explicitly stated for an arbitrary
subcategory $R'$ and every derivator, identifies precomposition by
$\gamma$ on cocontinuous morphisms into $\D$ with (B.3.2).
Precomposition by $i$ supplies its left adjoint. The same construction
with the opposite derivator and opposite subcategory supplies its
right adjoint. The internal-Hom argument of Theorem~3.13 gives these
as adjunctions of derivators, because the argument just given applies
to every shift. All the categories and sets of arrows used in the
localization are $\U$-small; the simplicial-presheaf categories are
taken in the ambient larger universe. This uses the universal
property of the auxiliary model derivator, not a presentation of
$\D$ by a model category.
There is no hidden strongness hypothesis in this universal property:
the cited foundational paper defines a two-sided derivator using
Der~1--4, and Proposition~3.12 makes it eligible for Corollary~3.26
on every small shape
\citep[Definitions~0.1--0.2, Proposition~3.12, and
Corollary~3.26]{Cisinski2008}.

We use the right adjoint of (B.3.2), denoted $Q_{T,R}$. Suppose that
$\pi:T\to S$ is a Grothendieck opfibration and let $R$ consist of
its vertical arrows, namely those sent to identities. Write
$\D(T)^{\mathrm{vert}}=\D(T;R)$ and $Q_T=Q_{T,R}$. Its objects
are exactly the diagrams locally constant on every fiber. For the
fiber inclusion $i_s:T_s\hookrightarrow T$, the canonical mate is
an isomorphism
\[
 i_s^*Q_T\ \cong\ Q_{T_s}i_s^*.
 \tag{B.3.4}
\]
We verify the required base change rather than assume descent.
The left Kan extension $i_{s!}$ sends locally constant objects on
$T_s$ to vertically locally constant objects on $T$. Indeed, at
$x\in T_t$, the comma category $x\backslash T_s$ is the disjoint
union, over $\alpha:t\to s$, of categories with initial object
the chosen cocartesian lift $x\to\alpha_!x$. Der~4, in presheaf
variance, therefore gives
\[
 x^*i_{s!}F\ \cong\
 \coprod_{\alpha\in\Hom_S(t,s)}(\alpha_!x)^*F.
 \tag{B.3.5}
\]
For a vertical arrow $x\to x'$, the comparison in (B.3.5) is the
coproduct of the comparisons induced by the vertical arrows
$\alpha_!x\to\alpha_!x'$. They are isomorphisms when $F$ is
locally constant. The assertion follows, including empty comma
categories and infinite coproducts. The equality of the two
composites ``include, then $i_{s!}$'' and ``$i_{s!}$ on the local
subcategories, then include'' now gives (B.3.4) by taking right
adjoints. Consequently (B.3.4) is compatible with the units and
counits, and the argument also proves it at every shift.

Consider a map $v:T\to T'$ of Grothendieck opfibrations over $S$
preserving cocartesian arrows. For $x'\in T'_s$, the inclusion
\[
 T_s/x'\hookrightarrow T/x'
 \tag{B.3.6}
\]
has a left adjoint. It sends $(x,\beta:v(x)\to x')$, with
$\pi(x)=t$ and $\pi'(\beta)=\alpha:t\to s$, to the object
$(\alpha_!x,\overline\beta)$ obtained by factoring $\beta$
through the cocartesian transport of $v(x)$. The universal property
of a cocartesian arrow gives the adjunction. If $r\dashv i$ is
this adjunction, presheaf variance gives $r_*\cong i^*$; composing
right Kan extensions to the point shows that restriction along
(B.3.6) computes the same homotopy limit. Der~4 therefore gives the
canonical isomorphism
\[
 i_s'^*v_*\ \cong\ v_{s*}i_s^*.
 \tag{B.3.7}
\]
This is the usual base change for a map of opfibrations, proved here
with the variance required for $\D$.

Suppose now that every $v_s^*:\LC_\D(T'_s)\to\LC_\D(T_s)$
is an equivalence. On vertically locally constant diagrams we have
an adjunction
\[
 v^*:\D(T')^{\mathrm{vert}}
 \rightleftarrows\D(T)^{\mathrm{vert}}:Q_{T'}v_*.
 \tag{B.3.8}
\]
By (B.3.4) and (B.3.7), its unit and counit restrict on every fiber
to those of
$v_s^*\dashv Q_{T'_s}v_{s*}$. Those units and counits are
isomorphisms by hypothesis. Every point lies in a fiber, so Der~2
makes the unit and counit of (B.3.8) isomorphisms. Thus (B.3.8) is
an equivalence.

If the fiber restrictions are only fully faithful, the same argument
applied to the unit alone proves that $v^*$ in (B.3.8) is fully
faithful. Its restriction to the full subcategories of fully locally
constant diagrams is therefore fully faithful as well. This proves
the asserted fiberwise full-faithfulness conclusion without an
essential-surjectivity or reflection argument.

It remains to restrict this equivalence to diagrams locally constant
on all arrows. Restriction $v^*$ preserves this condition. It also
reflects it among vertically locally constant diagrams. To see this,
let $N\in\D(T')^{\mathrm{vert}}$ and suppose $v^*N$ is locally
constant. For $\alpha:s\to t$, the cocartesian lifts define a
natural transformation between the corresponding fiber inclusions,
and hence a morphism
\[
 \alpha_!^*(i_t'^*N)\longrightarrow i_s'^*N
 \quad\hbox{in }\LC_\D(T'_s).
 \tag{B.3.9}
\]
Its image under $v_s^*$ is an isomorphism, since it is the transport
comparison for $v^*N$. The equivalence $v_s^*$ is conservative, so
(B.3.9) is an isomorphism. Every arrow of $T'$ factors as a
cocartesian arrow followed by a vertical arrow. Equation (B.3.9)
and vertical local constancy therefore show that $N$ is locally
constant. Essential surjectivity in (B.3.8), followed by this
reflection, gives essential surjectivity on fully locally constant
diagrams; full faithfulness is inherited from (B.3.8). We have proved
that $v^*:\LC_\D(T')\to\LC_\D(T)$ is an equivalence.

Finally replace $p:A\to S$ by the Grothendieck opfibration
\[
 \widetilde A=\int_{s\in S}(A/s)\longrightarrow S.
 \tag{B.3.10}
\]
An object is $(a,s,\alpha:p(a)\to s)$, and a morphism is a pair
$(h:a\to a',k:s\to s')$ with $k\alpha=\alpha'p(h)$.
Transport along $k$ replaces $\alpha$ by $k\alpha$. The forgetful
functor $r_A:\widetilde A\to A$ has the left adjoint
$a\mapsto(a,p(a),\id)$, with $r_Aj_A=\id$ and
$j_Ar_A\Rightarrow\id$. The latter natural transformation induces
an isomorphism on every locally constant object, by Der~2.
Consequently $r_A^*$ and $j_A^*$ are mutually inverse equivalences
on locally constant categories. The same holds for $B$. The functor
$u$ induces a transport-preserving map
$\widetilde u:\widetilde A\to\widetilde B$ whose fiber at $s$
is precisely $u/s$. Apply (B.3.8)--(B.3.9) to $\widetilde u$ and
then use the forgetful equivalences. For full faithfulness use the
unit-only version of (B.3.8) just proved. This proves both relative
Theorem A assertions.

The unit, counit, and base-change proof is the ordinary-derivator
specialization of the proof of
\citet[Main Theorem~2.5.5(2), pp.~52--54]{EF_Hoermann2015}.
The stable compact-generation assumptions in that fibered theorem
provide its partial Cartesian projectors. In the present
nonfibered situation, (B.3.2)--(B.3.3) provide those projectors for
every derivator, so those assumptions are unnecessary. Weak
saturation was proved in A.4. A category with a terminal object is
$\Wlc$-aspherical by the same adjoint-section argument as for
$r_A$, using its projection to the point and the terminal section.
Together with the relative theorem just proved, these are precisely
the three fundamental-localizer axioms.
\end{proof}
